% =====================================================================
\chapter{Les technologies utilisées}
\label{ch:technos}

Chaque technologie est justifiée par ce qu'elle apporte \emph{à ce projet},
avec le fichier qui la met en œuvre. Aucune définition générale.

% =====================================================================
\section{Backend}

\subsection{Python}
\textbf{Où} : tout le backend. Le README annonce Python 3.13.\\
\textbf{Pourquoi ici} : c'est la seule langue où cohabitent naturellement le SDK
MCP officiel, un client SSH (\texttt{paramiko}), un ORM mature (SQLAlchemy) et
les bibliothèques d'embeddings (\texttt{sentence-transformers}). Le projet
utilise aussi des constructions récentes du langage : \texttt{str | None},
\texttt{asyncio.to\_thread}, \texttt{@dataclass}, \texttt{match} implicite via
les dictionnaires de dispatch.

\subsection{FastAPI (+ Starlette, Uvicorn)}
\textbf{Où} : \file{app/main.py}, les 11 routeurs.\\
\textbf{Pourquoi ici} : trois raisons visibles dans le code.
\begin{itemize}
  \item \textbf{L'injection de dépendances} rend le contrôle d'accès
  déclaratif : \texttt{current\_user: User = Depends(require\_role(...))} suffit
  à protéger une route.
  \item \textbf{Le mélange \texttt{def} / \texttt{async def}} est exploité
  volontairement. Les routes qui font du SSH bloquant restent en \texttt{def}
  pour partir dans le threadpool (\file{dashboard.py}, \file{reports.py}) ; les
  routes qui attendent l'agent sont en \texttt{async def} (\file{storage.py},
  \route{/investigate}, \route{/confirm}). Les commentaires l'expliquent
  fichier par fichier.
  \item \textbf{Le \fn{lifespan}} héberge les trois tâches de fond avec
  annulation propre.
\end{itemize}

\subsection{Pydantic v2 et pydantic-settings}
\textbf{Où} : \file{app/schemas/}, \file{app/core/config.py}.\\
\textbf{Pourquoi ici} : la validation d'entrée est la première barrière.
\cls{ServiceRequestCreate} impose \texttt{title} entre 1 et 200 caractères ;
\cls{LoginRequest} valide le format d'e-mail (via \texttt{email-validator}) ;
\cls{ChatRequest} accepte un \texttt{conversation\_id} optionnel typé.
Côté configuration, \cls{Settings} \textbf{refuse de démarrer} si
\texttt{database\_url}, \texttt{jwt\_secret\_key} ou \texttt{groq\_api\_key}
manquent --- une panne au démarrage plutôt qu'à la première requête.

\begin{infobox}
Les schémas séparent volontairement la liste du détail : \cls{IncidentTicketOut}
pour la liste, \cls{IncidentTicketDetail} pour la fiche. Le docstring justifie :
« y joindre notes et actions ferait autant de requêtes supplémentaires que de
lignes affichées, pour des données que la liste ne montre pas. »
\end{infobox}

\subsection{SQLAlchemy (mode synchrone)}
\textbf{Où} : \file{app/models/}, \file{app/database/session.py}.\\
\textbf{Pourquoi \emph{synchrone}} : c'est un choix, et il est documenté. Les
docstrings de \file{incident\_service.py}, \file{mcp\_audit.py} et
\file{actions.py} portent tous la trace d'un correctif : une première version
avait été écrite pour \texttt{AsyncSession}, avec des \texttt{await
db.commit()} qui échouaient (« sur une session sync, \fn{db.commit()} renvoie
\texttt{None} et \texttt{await None} lève un \texttt{TypeError} »). Le projet a
convergé vers le synchrone, et FastAPI compense en exécutant les routes
\texttt{def} dans un threadpool.

\subsection{Alembic}
\textbf{Où} : \file{backend/alembic/}, 14 révisions.\\
\textbf{Pourquoi ici} : le schéma a beaucoup évolué (ajout du chat, des
conversations multiples, des demandes, des pièces jointes, du lien
investigation). Une migration comme \file{b6d0f3a91c47} fait ce qu'un
\texttt{create\_all} ne saurait pas faire : ajouter un index unique
\emph{partiel} PostgreSQL (\texttt{postgresql\_where}).

\subsection{PostgreSQL 16}
\textbf{Où} : \file{docker-compose.yml}, tous les modèles.\\
\textbf{Pourquoi ici} : JSONB interrogeable (\texttt{mcp\_data[key].astext}),
UUID natif, types ENUM, et surtout \texttt{date\_bin} --- fonction PostgreSQL
14+ sans laquelle le graphique de performance devrait être agrégé en Python.

\subsection{python-jose et passlib[bcrypt]}
\textbf{Où} : \file{app/core/security.py}.\\
\textbf{Pourquoi ici} : JWT HS256 pour une API sans état, et bcrypt pour ne
jamais stocker un mot de passe en clair. Le jeton porte le rôle, mais le contrôle
d'accès relit le rôle en base --- un changement de rôle prend effet
immédiatement.

\subsection{MCP (SDK officiel, version 2.0.0)}
\textbf{Où} : \file{app/mcp/storage/}.\\
\textbf{Pourquoi ici} : c'est le protocole qui permet d'exposer des capacités à
un modèle \emph{sans coder à la main} un catalogue de fonctions, un schéma JSON
par paramètre et un routeur d'appels. Le décorateur \texttt{mcp.tool()} dérive
tout des annotations de type et des docstrings. Le transport \texttt{stdio}
isole en outre les outils dans un processus séparé.

\subsection{Paramiko}
\textbf{Où} : \file{app/mcp/storage/ssh.py}.\\
\textbf{Pourquoi ici} : parce que le projet a fait le choix de \textbf{ne rien
simuler}. Il n'existe pas d'API REST sur une VM Ubuntu pour lire \texttt{df} ou
\texttt{lvs} : SSH est le protocole d'administration réel. Authentification par
clé, \texttt{AutoAddPolicy}, timeouts de 5\,s (connexion) et 15\,s (commande).

\subsection{Groq (SDK \texttt{groq})}
\textbf{Où} : \file{app/llm/groq\_client.py}.\\
\textbf{Pourquoi ici} : deux raisons opérationnelles. La \textbf{latence} ---
un ingénieur qui investigue un incident n'attend pas ; et le \textbf{tool
calling natif}, compatible avec le format OpenAI, ce qui permet de convertir le
catalogue MCP en \texttt{tools} d'une seule fonction. Le client essaie
successivement plusieurs modèles et retient le premier qui répond.

\subsection{ChromaDB}
\textbf{Où} : \file{app/rag/vectorstore.py}, conteneur \texttt{chromadb/chroma}.\\
\textbf{Pourquoi ici} : la base de connaissances est petite (quatre documents
d'origine plus les dépôts du Data Hub). Chroma s'installe en un conteneur,
accepte des embeddings fournis de l'extérieur, et supporte \texttt{upsert} ---
ce qui rend l'indexation relançable et la suppression d'un document possible.

\subsection{sentence-transformers}
\textbf{Où} : \file{app/rag/embedding\_service.py}.\\
\textbf{Pourquoi ce modèle précis} : le fichier l'explique. « la base de
connaissances Storage est en anglais (NetApp, Veeam, TrueNAS, AWS) alors que les
ingénieurs posent leurs questions en français \ldots{} Un modèle anglophone comme
\texttt{bge-small-en-v1.5} place ces deux langues dans des régions éloignées de
l'espace vectoriel : la question ne retrouve pas le passage qui y répond. »
D'où \texttt{intfloat/multilingual-e5-small}, avec ses préfixes asymétriques.

\subsection{ReportLab}
\textbf{Où} : \file{app/services/report\_service.py}.\\
\textbf{Pourquoi ReportLab et pas WeasyPrint} --- le docstring tranche :
« WeasyPrint s'appuie sur les bibliothèques GTK (libgobject, pango, cairo) qui
ne sont pas installées avec pip sous Windows, et son import échoue sur le poste
de développement. »

\begin{warnbox}
\textbf{Détail à corriger un jour.} \texttt{reportlab} figure dans
\file{requirements.lock.txt} mais \textbf{pas} dans \file{requirements.txt}.
Une installation faite depuis le seul \file{requirements.txt} ne pourrait pas
générer de rapport. C'est une observation, pas une supposition : la ligne est
absente.
\end{warnbox}

\subsection{pypdf}
\textbf{Où} : extraction du texte des PDF déposés au Data Hub.\\
\textbf{Pourquoi ici} : sans couche texte, pas d'indexation --- et le service le
dit à l'utilisateur en lui suggérant un OCR.

% =====================================================================
\section{Frontend}

\begin{longtable}{@{}p{3.2cm}p{2.0cm}p{9.5cm}@{}}
\toprule
\textbf{Technologie} & \textbf{Version} & \textbf{Pourquoi dans ce projet}\\
\midrule
\endhead
React & 18.3 & Interface fortement dynamique : un fil de chat qui grandit, des propositions d'action qui changent d'état après un clic, un drawer d'incident. L'état local suffit, aucun Redux n'est nécessaire\\
TypeScript & 5.7 & \file{lib/types.ts} (259 lignes) reproduit les schémas Pydantic. Renommer un champ côté backend casse la compilation côté frontend --- c'est le seul contrat automatique entre les deux moitiés du projet\\
Vite & 5.4 & Rechargement instantané, et surtout le \textbf{proxy} \route{/api} \arrowto{} \texttt{127.0.0.1:8000} qui supprime tout problème CORS en développement\\
TailwindCSS & 3.4 & Cohérence visuelle sans feuille de style séparée. Une palette \texttt{brand} / \texttt{ink} / \texttt{accent} / \texttt{danger} est réutilisée jusque dans les rapports PDF\\
react-router-dom & 6.30 & Routes imbriquées : \cls{RequireAuth} est une route parente qui rend \cls{AppLayout}, les pages sont ses enfants\\
recharts & 2.15 & Un seul graphique (\cls{PerformanceChart}), mais avec trois séries et trois plages temporelles\\
lucide-react & 0.475 & Icônes cohérentes dans la sidebar, les badges et les boutons\\
clsx + tailwind-merge & --- & \fn{cn()} : composer des classes conditionnelles sans conflit Tailwind\\
\bottomrule
\end{longtable}

% =====================================================================
\section{Infrastructure et outillage}

\subsection{Docker et Docker Compose}
\textbf{Où} : \file{docker-compose.yml}.\\
\textbf{Ce qui est réellement conteneurisé} : \texttt{postgres:16} (avec
\texttt{healthcheck} \texttt{pg\_isready} et volume nommé) et
\texttt{chromadb/chroma:latest} (port 8001 exposé, volume nommé).\\
\textbf{Pourquoi ici} : ces deux services sont pénibles à installer nativement,
surtout sous Windows. Le backend et le frontend, eux, tournent en local.

\subsection{SSH}
\textbf{Pourquoi ici} : c'est le protocole d'administration de la VM. Le projet
n'installe aucun agent sur la machine cible : il exécute les mêmes commandes
qu'un ingénieur taperait.

\subsection{Technologies Linux mobilisées sur la VM}

\begin{longtable}{@{}p{3.4cm}p{11.4cm}@{}}
\toprule
\textbf{Technologie} & \textbf{Rôle dans le projet}\\
\midrule
\endhead
LVM & Volumes logiques, snapshots (\texttt{lvcreate -{}-snapshot}), restauration (\texttt{lvconvert -{}-merge}). Lu par \texttt{lvs}\\
systemd & Les jobs de sauvegarde \emph{sont} des timers systemd. \texttt{list-timers}, \texttt{show}, \texttt{start}, les drop-ins et \texttt{revert}\\
journald & \fn{get\_backup\_logs} lit \texttt{journalctl -o short-iso}\\
\texttt{/proc/diskstats} & Source unique des métriques de performance\\
\texttt{df} & Source unique de la capacité\\
restic & Le logiciel de sauvegarde réel ; \fn{repair\_backup\_repository} sait lire l'état d'un dépôt\\
MinIO & Cible objet S3 du laboratoire, avec Object Lock\\
NFS / Samba & Partages du laboratoire NAS\\
\bottomrule
\end{longtable}

\begin{keybox}
\textbf{L'argument de fond, à retenir pour la soutenance.} Le choix
\texttt{systemd} + \texttt{restic} + \texttt{LVM} n'est pas un pis-aller : ce
sont les briques que l'outil interroge \emph{réellement}. Chaque relevé affiché
correspond à une commande qu'un jury pourrait rejouer lui-même en SSH sur la
VM.
\end{keybox}

% =====================================================================
\section{Ce qui est déclaré mais pas utilisé}

Un point d'honnêteté qui vaut mieux qu'une question embarrassante.

\begin{longtable}{@{}p{3.6cm}p{11.2cm}@{}}
\toprule
\textbf{Élément} & \textbf{Statut réel}\\
\midrule
\endhead
Anthropic Claude & Clé déclarée dans \cls{Settings}, \file{llm\_router.py} lève \texttt{NotImplementedError}. \textbf{Non branché}\\
\file{requirements.txt} & Contient un très grand nombre de paquets sans rapport (Jupyter, Streamlit, MLflow, XGBoost, OpenCV, Flask, Plotly\ldots), vestiges d'un environnement de travail global. La liste réellement nécessaire est bien plus courte\\
WeasyPrint & Présent dans les dépendances, remplacé par ReportLab pour la raison ci-dessus\\
\file{base\_persona.py}, \file{llm\_router.py}, \file{rbac\_middleware.py} & Code valide, importé par personne\\
\texttt{DESTRUCTIVE\_TOOLS}, \fn{risk\_for}, \fn{invoke} & Définis, jamais appelés\\
\tbl{request\_messages} & Table conservée, plus alimentée\\
Persona Linux & Le README l'annonce, et un commit récent (\texttt{df219f23}) porte « retrait de la persona Linux ». Seul le persona Storage existe dans le code\\
\bottomrule
\end{longtable}
